\documentclass[10pt,twocolumn]{article}
\usepackage[margin=0.75in]{geometry}
\usepackage{times}
\usepackage{amsmath, amssymb}
\usepackage{graphicx}
\usepackage{booktabs}
\usepackage{titlesec}
\usepackage{caption}
\usepackage{hyperref}

\titleformat{\section}{\bfseries}{\thesection.}{0.5em}{}
\titleformat{\subsection}{\bfseries}{\thesubsection}{0.5em}{}

\title{\textbf{Deep Learning for Stock Price Prediction:\\LSTM vs Transformer with Sentiment Analysis} \\[0.5em]
\large CS 7643 - Deep Learning}

\author{Agha A Raza, Konam G Kenneth \\
Georgia Institute of Technology \\
\texttt{ragha3@gatech.edu, kkonam3@gatech.edu}}

\date{}

\begin{document}

\maketitle

\begin{abstract}
\small
Stock price prediction remains challenging due to market noise and non-stationarity. We investigate whether deep learning models enhanced with sentiment analysis can improve next-day price direction prediction. We implement and train LSTM and Transformer architectures on four tech stocks (AAPL, AMZN, NVDA, TSLA) using technical indicators combined with StockTwits sentiment (2020-2022 training, 2026 validation). Key findings: (1) Both models perform nearly identically (LSTM: 43.57\%, Transformer: 44.00\%), with Transformer having minimal advantage (+0.43\%); (2) severe class imbalance causes prediction collapse, with models predicting predominantly one class; (3) both models perform below random baseline (50\%) on 2026 test data; (4) validation-to-test degradation of 13-15 percentage points indicates poor generalization; (5) limited training data (~415 samples) insufficient for either architecture. Our analysis demonstrates that financial time series prediction is significantly more challenging than initially apparent, requiring careful attention to class imbalance, data quantity, and evaluation methodology.
\end{abstract}

\section{Introduction}

\textbf{What are we trying to do?} Predict whether a stock's price will go up or down tomorrow, using both historical price data and investor sentiment from social media.

\textbf{Why is this hard?} Stock markets are noisy, non-stationary (patterns change over time), and influenced by countless factors. Traditional technical analysis uses only price/volume data, missing important sentiment signals from social media.

\textbf{Who cares?} Accurate price prediction could improve trading strategies and risk assessment. More broadly, understanding how deep learning handles financial time series has applications in economics and algorithmic trading.

\textbf{Our approach:} We train deep learning models (LSTM and Transformer) to learn temporal patterns from sequences of financial and sentiment features. We compare architectures, analyze failure modes, and validate generalization on 2026 out-of-sample data.

\textbf{Key contributions:}
\begin{itemize}
\item Rigorous comparison of LSTM vs Transformer on financial time series
\item Analysis of class imbalance and prediction collapse in stock prediction
\item Demonstration of poor generalization across market regimes
\item Evidence that limited data (~415 samples) insufficient for deep learning
\item Honest reporting of below-baseline results with methodological insights
\end{itemize}

\section{Related Work}

\textbf{Recurrent architectures} (LSTM, GRU) are popular for financial prediction due to their ability to capture temporal dependencies. \textbf{Attention mechanisms} and Transformers have shown promise in NLP and are increasingly applied to time series. \textbf{Sentiment analysis} from social media (Twitter, StockTwits) has been shown to provide predictive signals for stock movements. Our work extends prior research by: (1) implementing Transformer from scratch for financial prediction, (2) rigorous train/validation/test methodology, (3) comprehensive analysis of failure modes including class imbalance and prediction collapse, and (4) honest reporting of negative results.

\section{Data}

\subsection{Dataset Description}

We collected daily data for four large-cap technology stocks: AAPL, AMZN, NVDA, and TSLA.

\textbf{Training}: 2020-2022 (755 trading days, 80/20 train/val split)\\
\textbf{Validation}: 2020-2022 (20\% holdout, ~104 samples per ticker)\\
\textbf{Test}: 2026 (53-58 trading days, out-of-sample)\\
\textbf{Total}: ~2,600 stock-days across 4 tickers

\subsection{Data Collection}

\textbf{Market Data:} Yahoo Finance API provides adjusted closing prices, volume, and OHLC data. We compute 17 technical indicators including RSI, MACD, Bollinger Bands, and rolling returns.

\textbf{StockTwits Sentiment:} Public API provides social media messages with user-labeled sentiment (bullish/bearish). We score messages using RoBERTa (\texttt{cardiffnlp/twitter-roberta-base-sentiment-latest}) and aggregate daily: message count, sentiment mean/std, and sum statistics. Coverage: 755 days (100\%).

\subsection{Features and Target}

\textbf{Input Features (31 total):}
\begin{itemize}
\item Technical (17): Price, returns, volatility, RSI, MACD, Bollinger Bands, volume
\item Sentiment (14): StockTwits metrics (count, mean, std, sum)
\end{itemize}

\textbf{Target:} Binary classification - UP (1) if $P_{t+1} > P_t$, DOWN (0) otherwise.

\textbf{Sequence Construction:} Rolling windows of 5 days (31 features × 5 timesteps) predict next-day direction.

\subsection{Data Limitations}

\begin{itemize}
\item Only 4 tech stocks (may not generalize to other sectors)
\item Limited training data (~415 samples after 80/20 split)
\item Survivorship bias (only large-cap stocks)
\item Market regime: trained on bull market (2020-2022), tested on 2026
\end{itemize}

\section{Methodology}

\subsection{LSTM Architecture}

\textbf{Architecture:}
\begin{itemize}
\item Input: 31 features × 5 timesteps
\item LSTM: 2 layers, 128 hidden units, 20\% dropout
\item Dense: 64 units, ReLU activation, 20\% dropout
\item Output: 1 unit, Sigmoid activation
\end{itemize}

\textbf{Parameters:} ~200K trainable

\subsection{Transformer Architecture}

\textbf{Architecture:}
\begin{itemize}
\item Input projection: 31 → 64 dimensions
\item Positional encoding: Sinusoidal
\item Transformer encoder: 2 layers, 4 attention heads
\item Feed-forward: 128 dimensions per layer
\item Global average pooling
\item Classification head: 64 → 32 → 1
\end{itemize}

\textbf{Parameters:} ~150K trainable

\subsection{Training Details}

\textbf{Train/Val/Test Split:}
\begin{itemize}
\item Split 2020-2022 into train (80\%, ~415 samples) and validation (20\%, ~104 samples)
\item Select best model based on validation accuracy
\item Evaluate on 2026 test data only once at the end
\end{itemize}

\textbf{Loss Function:} Binary cross-entropy

\textbf{Optimizer:} Adam with learning rate 0.001

\textbf{Regularization:}
\begin{itemize}
\item Dropout: 20\% in LSTM, 10\% in Transformer
\item Early stopping: patience=10 epochs on validation loss
\end{itemize}

\textbf{Hyperparameters:}
\begin{itemize}
\item Sequence length: 5 days
\item Batch size: 32
\item Max epochs: 50 (early stopping typically at 10-20)
\end{itemize}

\section{Results}

\subsection{Overall Performance}

\begin{table}[h]
\centering
\small
\begin{tabular}{lcc}
\toprule
Model & Validation & Test (2026) \\
\midrule
LSTM & 58.63\% & 43.57\% \\
Transformer & 57.43\% & 44.00\% \\
\midrule
Random Baseline & 50.00\% & 50.00\% \\
\bottomrule
\end{tabular}
\caption{Overall model performance. Both models perform below random baseline on 2026 test data.}
\end{table}

\textbf{Key Finding:} Both models perform nearly identically, with Transformer having minimal advantage (+0.43\%). More critically, both perform below 50\% random baseline on test data.

\subsection{Per-Ticker Results}

\begin{table}[h]
\centering
\small
\begin{tabular}{lcccc}
\toprule
Ticker & \multicolumn{2}{c}{LSTM} & \multicolumn{2}{c}{Transformer} \\
\cmidrule(lr){2-3} \cmidrule(lr){4-5}
 & Val & Test & Val & Test \\
\midrule
AAPL & 59.62\% & 48.28\% & 58.65\% & 48.28\% \\
AMZN & 53.33\% & 41.38\% & 52.38\% & 41.38\% \\
NVDA & 62.50\% & 41.51\% & 62.50\% & 41.51\% \\
TSLA & 59.05\% & 43.10\% & 56.19\% & 44.83\% \\
\midrule
\textbf{Avg} & \textbf{58.63\%} & \textbf{43.57\%} & \textbf{57.43\%} & \textbf{44.00\%} \\
\bottomrule
\end{tabular}
\caption{Per-ticker performance. Models perform identically on AAPL, AMZN, and NVDA. Only TSLA shows difference (+1.73\% for Transformer).}
\end{table}

\textbf{Analysis:} 
\begin{itemize}
\item AAPL, AMZN, NVDA: Identical performance (both models predict same class)
\item TSLA: Transformer slightly better (+1.73\%)
\item All tickers: Massive validation-to-test degradation (-11\% to -21\%)
\end{itemize}

\subsection{Class Imbalance Analysis}

Both models suffer from severe class imbalance:

\textbf{AAPL, AMZN, NVDA (both models):}
\begin{itemize}
\item Down: 100\% recall, 41-48\% precision
\item Up: 0\% recall, 0\% precision
\item \textbf{Models always predict "Down"}
\end{itemize}

\textbf{TSLA:}
\begin{itemize}
\item LSTM: Down 88\% recall, Up 6\% recall
\item Transformer: Down 100\% recall, Up 0\% recall
\item \textbf{Models predominantly predict "Down"}
\end{itemize}

\textbf{Why this happens:}
\begin{itemize}
\item Test period has class imbalance (more "Down" days)
\item Models learn to predict majority class
\item Achieves reasonable accuracy but no predictive value
\end{itemize}

\subsection{Generalization Analysis}

\begin{table}[h]
\centering
\small
\begin{tabular}{lcc}
\toprule
Ticker & LSTM Deg. & Trans. Deg. \\
\midrule
AAPL & -11.34 pp & -10.37 pp \\
AMZN & -11.95 pp & -11.00 pp \\
NVDA & -20.99 pp & -20.99 pp \\
TSLA & -15.95 pp & -11.36 pp \\
\midrule
\textbf{Average} & \textbf{-15.06 pp} & \textbf{-13.43 pp} \\
\bottomrule
\end{tabular}
\caption{Validation-to-test degradation. Both models show massive performance drop, indicating poor generalization to 2026 market conditions.}
\end{table}

\textbf{Analysis:} 13-15 percentage point degradation indicates:
\begin{itemize}
\item Market regime change between 2020-2022 and 2026
\item Models overfitting to training period patterns
\item Insufficient regularization or data diversity
\end{itemize}

\section{Discussion}

\subsection{Why Both Models Fail}

\textbf{1. Insufficient Training Data}
\begin{itemize}
\item Only ~415 training samples after 80/20 split
\item Deep learning typically requires 10K-100K samples
\item Both LSTM and Transformer underperform
\end{itemize}

\textbf{2. Severe Class Imbalance}
\begin{itemize}
\item Models exploit test set class distribution
\item Predict majority class to minimize loss
\item No real pattern learning occurs
\end{itemize}

\textbf{3. Market Regime Change}
\begin{itemize}
\item Training: 2020-2022 bull market
\item Testing: 2026 mixed conditions
\item Patterns don't transfer across regimes
\end{itemize}

\textbf{4. Limited Signal in Sentiment}
\begin{itemize}
\item Sentiment alone insufficient for prediction
\item Need additional features (news, fundamentals, macro)
\item Social media sentiment may be noisy
\end{itemize}

\subsection{Why Models Perform Identically}

For 3 out of 4 tickers, both models achieve identical accuracy because:
\begin{enumerate}
\item Both predict only the "Down" class
\item Accuracy equals proportion of "Down" samples in test set
\item No architectural advantage when prediction collapses
\end{enumerate}

Only TSLA shows difference because Transformer maintains slightly more balanced predictions.

\subsection{Comparison to Literature}

Prior work reports 52-58\% accuracy on stock prediction tasks. Our results (43-44\%) are below this range, likely due to:
\begin{itemize}
\item Stricter evaluation (proper train/val/test split)
\item Limited training data
\item Challenging test period (2026 market conditions)
\item No data leakage or look-ahead bias
\end{itemize}

\subsection{Limitations}

\begin{enumerate}
\item \textbf{Limited scope}: Only 4 tech stocks
\item \textbf{Insufficient data}: ~415 training samples too few
\item \textbf{Short horizon}: Next-day prediction only
\item \textbf{Class imbalance}: Not addressed in training
\item \textbf{No transaction costs}: Real trading has fees
\item \textbf{Single sentiment source}: Only StockTwits
\end{enumerate}

\section{Lessons Learned}

\subsection{Methodological Insights}

\textbf{1. Proper Train/Val/Test Split is Critical}
\begin{itemize}
\item Validation set must be used for model selection
\item Test set should be "locked box" until final evaluation
\item Prevents overfitting to test set patterns
\end{itemize}

\textbf{2. Class Imbalance Must Be Addressed}
\begin{itemize}
\item Use class weights in loss function
\item Try focal loss or resampling techniques
\item Monitor per-class metrics, not just accuracy
\end{itemize}

\textbf{3. Multiple Metrics Required}
\begin{itemize}
\item Accuracy alone insufficient
\item Check precision, recall, F1 per class
\item Monitor prediction distribution
\end{itemize}

\subsection{Practical Recommendations}

\textbf{For Future Work:}
\begin{enumerate}
\item \textbf{Collect more data}: Need 10-100x more samples
\item \textbf{Address class imbalance}: Use class weights or focal loss
\item \textbf{Try simpler baselines}: Logistic regression, random forest
\item \textbf{Add more features}: News, fundamentals, market-wide indicators
\item \textbf{Ensemble methods}: Combine multiple models
\item \textbf{Market regime detection}: Separate models for different conditions
\end{enumerate}

\subsection{When to Use LSTM vs Transformer}

Based on our results:

\textbf{Use LSTM when:}
\begin{itemize}
\item Limited training data (<10K samples)
\item Strong temporal dependencies
\item Need stable training
\end{itemize}

\textbf{Use Transformer when:}
\begin{itemize}
\item Large training data (>100K samples)
\item Need interpretability (attention weights)
\item Parallel processing beneficial
\end{itemize}

\textbf{Our case:} With ~415 samples, neither architecture has sufficient data. Performance is nearly identical.

\section{Conclusion}

We implemented and trained LSTM and Transformer architectures for stock price prediction with sentiment analysis. Our findings:

\textbf{Key Results:}
\begin{enumerate}
\item \textbf{Both models perform identically} (LSTM: 43.57\%, Transformer: 44.00\%)
\item \textbf{Below random baseline} (50\%) on 2026 test data
\item \textbf{Severe class imbalance} causes prediction collapse
\item \textbf{Poor generalization} (-13\% to -15\% degradation)
\item \textbf{Insufficient training data} (~415 samples too few)
\end{enumerate}

\textbf{Insights:}
\begin{itemize}
\item Financial time series prediction is harder than initially apparent
\item Deep learning requires significantly more data than available
\item Class imbalance must be addressed before meaningful progress
\item Proper validation methodology is critical
\item Negative results are valuable for the field
\end{itemize}

\textbf{Future Work:}
\begin{itemize}
\item Collect 10-100x more training data
\item Address class imbalance with focal loss or resampling
\item Try simpler baselines (logistic regression, random forest)
\item Add more features (news, fundamentals, macro indicators)
\item Explore ensemble methods and transfer learning
\end{itemize}

Our work demonstrates that rigorous evaluation methodology and honest reporting of negative results are essential for advancing deep learning research. The fact that both sophisticated architectures perform below baseline highlights the fundamental challenges in financial time series prediction and the importance of data quantity and quality.

\section*{Acknowledgments}

We thank the CS 7643 teaching staff for guidance throughout this project. StockTwits data collected via public API. RoBERTa sentiment model from HuggingFace.

\section*{Code Repository}

All code, trained models, and data available at:\\
\url{https://github.com/KKONAM/financial-sentiment-analysis-spring-2026/tree/transformer-model-2026}

\newpage

\section*{References}

\begin{enumerate}
\item Hochreiter, S., \& Schmidhuber, J. (1997). Long short-term memory. Neural computation, 9(8), 1735-1780.
\item Vaswani, A., et al. (2017). Attention is all you need. Advances in neural information processing systems, 30.
\item Goodfellow, I., Bengio, Y., \& Courville, A. (2016). Deep learning. MIT press.
\item Bollen, J., Mao, H., \& Zeng, X. (2011). Twitter mood predicts the stock market. Journal of computational science, 2(1), 1-8.
\item Kingma, D. P., \& Ba, J. (2014). Adam: A method for stochastic optimization. arXiv preprint arXiv:1412.6980.
\end{enumerate}

\newpage

\section*{Team Contributions}

\begin{table}[h]
\centering
\begin{tabular}{p{3cm}p{10cm}}
\toprule
\textbf{Team Member} & \textbf{Contributions} \\
\midrule
Agha A Raza & 
\begin{itemize}
\item Implemented Transformer architecture from scratch
\item Designed and conducted experiments
\item Analyzed class imbalance and prediction collapse
\item Wrote methodology and analysis sections
\item Created training visualizations
\end{itemize} \\
\midrule
Konam G Kenneth & 
\begin{itemize}
\item Implemented LSTM baseline architecture
\item Collected and processed StockTwits data with RoBERTa
\item Designed proper train/validation/test methodology
\item Conducted 2026 out-of-sample validation
\item Wrote data collection and results sections
\end{itemize} \\
\midrule
\textbf{Shared} & 
\begin{itemize}
\item Hyperparameter tuning and model training
\item Analysis of results and failure modes
\item Paper writing and editing
\item Code repository organization
\item Figure generation and visualization
\end{itemize} \\
\bottomrule
\end{tabular}
\end{table}

\textbf{Note:} Both team members contributed equally to the technical implementation, experimentation, and analysis.

\end{document}
